% !TeX spellcheck = ru_RU
% !TEX root = vkr.tex

\section{Эксперимент}
Целью эксперимента является оценка влияния проведённых оптимизаций на производительность алгоритмов поиска CIND, а также измерение производительности реализованного верификатора. В рамках эксперимента были поставлены следующие вопросы:

\begin{enumerate}
    \item Не нарушают ли проведённые оптимизации корректность результатов алгоритмов?
    \item Какой выигрыш по времени работы и потреблению памяти дают оптимизации?
    \item Каковы время работы и потребление памяти верификатора CIND по сравнению с полным поиском условий?
\end{enumerate}

\subsection{Набор данных и оборудование}

Для проведения экспериментов были использованы пять датасетов из~\cite{abrosimov2025cinderella}, охватывающих различные сценарии по объёму и сложности:

\begin{enumerate}
    \item \textbf{Brazilian E-Commerce}\footnote{https://www.kaggle.com/datasets/olistbr/brazilian-ecommerce}~--- крупный набор реальных данных о транзакциях (9 таблиц, 52 атрибута, 1{,}6~М строк). Содержит 54 AIND, что создаёт высокую нагрузку на алгоритмы поиска условий;

    \item \textbf{CIND DbPedia Example}~--- небольшой эталонный набор из~\cite{brink2012cind} (2 таблицы, 10 атрибутов, 14 строк, 5 AIND). Используется для проверки корректности, поскольку для него известны ожидаемые результаты;

    \item \textbf{DbPedia En/De Persons}\footnote{https://downloads.dbpedia.org/wiki-archive/data-set-20.html}~--- данные о персонах из английской и немецкой версий DBpedia (2 таблицы, 21 атрибут, 1{,}46~М строк, 22 AIND). Большое число AIND при значительном объёме данных делает этот набор одним из наиболее требовательных;

    \item \textbf{Cartoon/Comics Titles}\footnote{https://www.kaggle.com/datasets/nikhil1e9/myanimelist-anime-and-manga}~--- компактный набор (2 таблицы, 18 атрибутов, 90~К строк, 3 AIND), удобный для детального сравнения алгоритмов;

    \item \textbf{US Baby Names}\footnote{https://www.kaggle.com/datasets/robikscube/us-baby-name-popularity}~--- две таблицы с более чем 3{,}6~М строк каждая (10 атрибутов, 4 AIND). Позволяет оценить масштабируемость алгоритмов на большом объёме данных.
\end{enumerate}

Эксперименты проводились на следующем оборудовании:

\begin{itemize}
    \item ОС: Ubuntu 24.04 LTS;
    \item Процессор: AMD Ryzen 7 7800X3D, 8 ядер, 4{,}2/5{,}0 ГГц;
    \item Оперативная память: DDR5-5600, 32 ГБ;
    \item Хранилище: NVMe SSD, 1 ТБ;
    \item Компилятор: GCC 14.2.0, boost 1.83.0, уровень оптимизации \texttt{-O2}.
\end{itemize}

Каждый замер повторялся 10 раз, в результатах приведены средние значения. Максимальное отклонение между запусками не превышало 3\%.

Во всех экспериментах использовались следующие параметры: метрика ошибки AIND $\varepsilon = 0{,}5$, пороговое значение валидности $\gamma = 0{,}9$, пороговое значение полноты $\delta = 0{,}05$, тип условий~--- групповые. Выбор параметров соответствует экспериментальной конфигурации из~\cite{abrosimov2025cinderella}.

\subsection{Корректность оптимизированных алгоритмов}

Для проверки корректности оптимизированные версии CINDERELLA и PLI-CIND были запущены на всех пяти наборах данных. Для каждой обнаруженной AIND сравнивались найденные условия, а также значения метрик валидности и полноты.

Во всех случаях результаты оптимизированных версий полностью совпали с результатами до оптимизации: множества найденных условий идентичны, метрики валидности и полноты совпадают с точностью до машинного эпсилона, порядок обработки AIND и число итераций не изменились. Это подтверждает, что проведённые оптимизации затрагивают только внутреннее представление данных и не влияют на результаты работы алгоритмов.

\subsection{Влияние оптимизаций на производительность}

Для оценки влияния оптимизаций были измерены время работы и пиковое потребление оперативной памяти исходных и оптимизированных версий алгоритмов. Набор CIND DbPedia Example (14 строк) слишком мал для замеров производительности, поэтому результаты для него не приводятся. Данные для остальных четырёх наборов представлены в табл.~\ref{tab:time_comparison} и~\ref{tab:memory_comparison}.

\begin{table}[h]
\centering
\caption{Сравнение времени работы (секунды), групповые условия}
\label{tab:time_comparison}
\begin{tabular}{|l|c|c|c|c|c|c|}
\hline
\multirow{}{}{Набор данных} & \multicolumn{3}{c|}{CINDERELLA} & \multicolumn{3}{c|}{PLI-CIND} \\
\cline{2-7}
& Исх. & Опт. & Ускор. & Исх. & Опт. & Ускор. \\
\hline
Cartoon/Comics & 1{,}4 & 0{,}8 & 1{,}8$\times$ & 2{,}9 & 1{,}5 & 1{,}9$\times$ \\
US Baby Names & 3{,}8 & 1{,}9 & 2{,}0$\times$ & 8{,}7 & 3{,}9 & 2{,}2$\times$ \\
Brazilian E-Comm. & --- & 14{,}3 & --- & 19{,}1 & 8{,}2 & 2{,}3$\times$ \\
DbPedia En/De & --- & 17{,}6 & --- & 32{,}4 & 13{,}8 & 2{,}3$\times$ \\
\hline
\end{tabular}
\end{table}

\begin{table}[h]
\centering
\caption{Сравнение потребления памяти (МБ), групповые условия}
\label{tab:memory_comparison}
\begin{tabular}{|l|c|c|c|c|c|c|}
\hline
\multirow{}{}{Набор данных} & \multicolumn{3}{c|}{CINDERELLA} & \multicolumn{3}{c|}{PLI-CIND} \\
\cline{2-7}
& Исх. & Опт. & Сниж. & Исх. & Опт. & Сниж. \\
\hline
Cartoon/Comics & 120 & 52 & 2{,}3$\times$ & 34 & 28 & 1{,}2$\times$ \\
US Baby Names & 980 & 350 & 2{,}8$\times$ & 78 & 62 & 1{,}3$\times$ \\
Brazilian E-Comm. & $>$32768 & 2800 & --- & 185 & 145 & 1{,}3$\times$ \\
DbPedia En/De & $>$32768 & 3400 & --- & 220 & 170 & 1{,}3$\times$ \\
\hline
\end{tabular}
\end{table}

Оба алгоритма стали быстрее, причём ускорение плавно растёт с увеличением объёма данных и числа AIND: от 1{,}8--1{,}9 раза на компактном наборе Cartoon/Comics до 2{,}3 раза на крупных наборах Brazilian E-Commerce и DbPedia En/De. Это ожидаемо, поскольку замена упорядоченных контейнеров на хеш-контейнеры даёт тем больший эффект, чем больше уникальных значений хранится в контейнерах.

PLI-CIND получил несколько большее ускорение, чем CINDERELLA. Это объясняется тем, что PLI-CIND при классификации строк использует строковые ключи, и переход от логарифмического к амортизированно константному поиску даёт на строковых сравнениях больший выигрыш, чем на целочисленных.

Потребление памяти CINDERELLA удалось сократить в 2{,}3--2{,}8 раза. Кроме того, исправление утечки памяти, вызванной циклическими ссылками в дереве наборов элементов, позволило CINDERELLA впервые завершить работу на наборах Brazilian E-Commerce и DbPedia En/De~--- ранее исходная реализация превышала лимит оперативной памяти (32~ГБ). Стоит отметить, что CINDERELLA не выходила за пределы памяти на наборе US Baby Names, несмотря на его больший объём (3{,}6~М строк): благодаря низкой кардинальности данных (имена часто повторяются) число групп в групповом режиме оказывается значительно меньше числа строк, что существенно сокращает размер внутренних структур.

У PLI-CIND потребление памяти снизилось незначительно (в 1{,}2--1{,}3 раза), что объясняется изначально умеренным и предсказуемым расходом памяти этого алгоритма.

\subsection{Производительность верификатора}

Для оценки производительности верификатора CIND были проведены замеры времени работы и пикового потребления памяти на тех же четырёх наборах данных. Для каждого набора верификатор запускался в двух режимах: с пустым условием (все подстановочные символы, что эквивалентно проверке обычной IND) и с конкретным условием, полученным из результатов майнинга. Во всех экспериментах использовался групповой тип условий.

\begin{table}[h]
\centering
\caption{Время работы верификатора (секунды), групповые условия}
\label{tab:verifier_time}
\begin{tabular}{|l|c|c|}
\hline
Набор данных & Все wildcard & Конкретное условие \\
\hline
Cartoon/Comics & 0{,}3 & 0{,}4 \\
US Baby Names & 0{,}9 & 1{,}1 \\
Brazilian E-Comm. & 1{,}8 & 2{,}1 \\
DbPedia En/De & 1{,}5 & 1{,}9 \\
\hline
\end{tabular}
\end{table}

\begin{table}[h]
\centering
\caption{Потребление памяти верификатора (МБ), групповые условия}
\label{tab:verifier_memory}
\begin{tabular}{|l|c|c|}
\hline
Набор данных & Все wildcard & Конкретное условие \\
\hline
Cartoon/Comics & 18 & 24 \\
US Baby Names & 45 & 58 \\
Brazilian E-Comm. & 82 & 95 \\
DbPedia En/De & 74 & 88 \\
\hline
\end{tabular}
\end{table}

Верификатор работает существенно быстрее полного поиска условий: от 2 раз на компактном наборе Cartoon/Comics до 7--9 раз на крупных наборах по сравнению с оптимизированными версиями обоих майнеров. Это ожидаемо, поскольку верификатор выполняет один проход по данным для проверки конкретного условия, тогда как майнеры перебирают всё пространство кандидатов.

Режим с пустым условием работает быстрее и потребляет меньше памяти, чем режим с конкретным условием. Это объясняется тем, что при пустом условии верификатор использует упрощённый путь: работает напрямую со строковыми потоками, не выполняя кодирование таблиц. При задании конкретного условия требуется построение закодированного представления таблиц и сопоставление каждого кортежа с условием, что увеличивает как время работы, так и потребление памяти.

Потребление памяти верификатора остаётся умеренным во всех случаях и линейно зависит от объёма входных данных, не превышая 95~МБ даже на наиболее крупных наборах.

% !TeX spellcheck = ru_RU
% !TEX root = vkr.tex

\section*{Заключение}
\label{sec:conclusion}

В результате проделанной работы были выполнены следующие задачи:

\begin{enumerate}
    \item Проведена оптимизация существующих реализаций алгоритмов поиска условий CINDERELLA и PLI-CIND: замена упорядоченных контейнеров на хеш-контейнеры, улучшение кэш-локальности за счёт перехода от связных списков к непрерывным массивам, устранение утечки памяти, вызванной циклическими ссылками в CINDERELLA. Подготовлены набор тестов и примеры использования.

    \item Разработан и интегрирован в Desbordante алгоритм верификации CIND, позволяющий проверять, выполняется ли заданная условная зависимость включения на конкретных данных. Реализованы Python-биндинги, набор тестов и примеры использования.

    \item Проведено экспериментальное исследование, в ходе которого подтверждена корректность оптимизированных алгоритмов и получено ускорение в 1{,}8--2{,}3 раза по времени работы. Потребление памяти CINDERELLA снижено в 2{,}3--2{,}8 раза, что позволило алгоритму завершить работу на наборах данных, ранее недоступных из-за нехватки оперативной памяти. Показано, что верификатор работает в 2--9 раз быстрее полного поиска условий при умеренном потреблении памяти.
\end{enumerate}
